\scanpage{59}
Les groupes d'automorphismes de ces MD-graphes sont
\figure{\begin{tikzpicture}[x=1cm,y=1cm]
\node at (-1.25,-1) {$(73)$};\draw (-.7,.6) .. controls (-.9,.6) and (-.8,0) .. (-.8,-.5) .. controls (-.8,-.7) and (-.9,-.8) .. (-1,-.8) .. controls (-.8,-.9) and (-.8,-1.1) .. (-.75,-1.5) -- (-.65,-2.4);
\draw (0,.35)--(1.5,.35);\fill (0,.35) circle(1.5pt);\fill (1.5,.35) circle(1.5pt);\node[above] at (0,.35) {$3$};\node[above] at (1.5,.35) {$\unsure{3|2}$};\node[below] at (0,.35) {$\alpha_0$};\node[right] at (2.4,.4) {$\Gamma=\unsure{\mathfrak S}_3\times\mathfrak S_2,\quad\Gamma^!=1,\quad\operatorname{Im}(\Gamma\to\mathfrak S_4)\simeq\mathfrak S_3$};
\draw (0,-.85) .. controls (2,.3) and (2.25,-1.3) .. (1,-.95) -- (0,-.85);\fill (0,-.85) circle(1.5pt);\node[above] at (0,-.85) {$2$};\node[below] at (0,-.85) {$(\alpha_0)$};\node[right] at (2.4,-.8) {$\Gamma\simeq\mathfrak S_2\times\mathfrak S_2,\quad\Gamma^!\simeq\mathfrak S_2,\quad\operatorname{Im}(\Gamma\to\mathfrak S_4)\simeq\underset{\substack{\Vert\\\mathfrak S_2\times\mathfrak S_2}}\Gamma$};
\draw (0,-2.2)--(1.7,-2.2);\fill (0,-2.2) circle(1.5pt);\fill (1.7,-2.2) circle(1.5pt);\node[above] at (1.7,-2.2) {$2$};\node[below] at (0,-2.2) {$g_{\alpha_0}=1$};\node[right] at (2.4,-2.1) {$\Gamma\simeq\mathfrak S_2,\quad\Gamma^!\simeq1\quad\deleted{\unsure{\operatorname{Im}(\Gamma\to\mathfrak S_4)=1}}$};
\end{tikzpicture}}

III) \uline{Dim. modulaire $2$}

Correspond par (67) aux deux cas suivants :

$1^\circ)$ Tous les $\alpha\in S$ sauf un seul \inserted{$\alpha_0$} sont de genre $0$, et de poids total $3$ (i.e. d'ordre \inserted{$\widetilde\mu_\alpha$} $\leq3$, et de poids marqué $3-\widetilde\mu_\alpha$), et $(\unsure{\alpha_0|g_{\alpha_0}},\widehat\nu_{\alpha_0})$ est d'un des deux types $(0,5)$, $(1,2)$ \deleted{$\unsure{(2,0)}$}, i.e. $\alpha_0$ est de genre $0$, d'ordre \inserted{$\widetilde\mu_{\alpha_0}$} $\leq5$ et de poids marqué $5-\widetilde\mu_{\alpha_0}$, \uline{ou} $\alpha_0$ de genre $1$, d'ordre \deleted{$\unsure{\leq2}$} \inserted{$\widetilde\mu_{\alpha_0}\in\{1,2\}$} et de poids marqué $2-\widetilde\mu_{\alpha_0}\in\{1,0\}$.
\margin{donc deux cas en réalité}

$2^\circ)$ Tous les $\unsure{\alpha=S|\alpha\in S}$ sauf deux exactement $\alpha_0$ et $\alpha_1$, sont de genre $0$, \inserted{et} de poids total $3$ (i.e. d'ordre $1\leq\widetilde\mu_\alpha\leq3$, et de poids marqué égal à $3-\widetilde\mu_\alpha$), $\alpha_0$ et $\alpha_1$ sont chacun \deleted{soit} de type $(\unsure{\alpha_i|g_i},\unsure{\widehat\mu_i|\widehat\nu_i})$ égal soit $(0,4)$, soit à $(1,1)$.
\margin{Cela fait donc trois cas, suivant que les deux types de dim mod $1$ qui interviennent sont tous deux $(0,4)$, ou l'un $(0,4)$, l'autre $(1,1)$, ou tous deux $(1,1)$.}

La \deleted{\unsure{schéma|variété}} \inserted{mult.} modulaire $\widehat M_G$ est alors isomorphe, suivant les cas, à l'une des suivantes :
\margin{58}
